% !TEX program = lualatex
\documentclass[11pt,a4paper]{article}

% ========= 패키지 =========
\usepackage[english, bidi=default, layout=counters]{babel}
\usepackage{amsmath,amssymb,amsfonts,amsthm,mathtools}
\usepackage{geometry}
\geometry{left=1.25cm,right=1.25cm,top=1.25cm,bottom=3.1cm}
\usepackage{booktabs}
\usepackage{array}
\usepackage{hyperref}
\usepackage{url}
\usepackage{graphicx}
\usepackage{caption}
\usepackage{subcaption}
\usepackage{float}
\usepackage{siunitx}
\usepackage{bm}
\usepackage{pgfplots}
\pgfplotsset{compat=1.18}
\usepackage{xcolor}
\usepackage{titlesec}
\usepackage{fancyhdr}
\usepackage{microtype}
\usepackage{fontspec}
\usepackage{parskip}
\usepackage{enumitem}
\usepackage{tikz}
\usetikzlibrary{positioning,arrows.meta}
\usepackage[numbers]{natbib}
\usepackage{xurl}

% ========= 한국어 폰트 설정 (Windows) =========
\usepackage{luatexja}
\usepackage{luatexja-fontspec}
\defaultfontfeatures{Ligatures=TeX}
\setmainfont{Times New Roman}[Script=Latin, Language=English]
\setsansfont{Malgun Gothic}[Script=CJK, Language=Korean]
\setmonofont{Courier New}[Script=Latin, Language=English]
\setmainjfont{Malgun Gothic}[Script=CJK, Language=Korean]
\setsansjfont{Malgun Gothic}[Script=CJK, Language=Korean]
\setmonojfont{Noto Sans Mono CJK KR}[Script=CJK, Language=Korean]

% ========= 언어 (한국어를 메인으로) =========
\babelprovide[import, main]{korean}
\babelprovide[import, onchar=ids fonts]{english}

% ========= 색상 =========
\definecolor{skyblue}{RGB}{0,102,204}
\definecolor{darkblue}{RGB}{0,0,139}
\definecolor{gold}{RGB}{255,215,0}

% ========= YUCT 매크로 =========
\newcommand{\Keff}{K_{\mathrm{eff}}}
\newcommand{\Dir}{\mathcal{D}}
\newcommand{\Mpl}{M_{\mathrm{Pl}}}
\newcommand{\Lpl}{l_{\mathrm{Pl}}}
\newcommand{\tP}{t_{\mathrm{P}}}
\newcommand{\Sodd}{S_{\mathrm{odd}}}
\newcommand{\Seven}{S_{\mathrm{even}}}
\newcommand{\kappac}{\kappa_{c}}
\newcommand{\betaf}{\beta}
\newcommand{\lagr}{\mathcal{L}}
\newcommand{\dint}{D\!+\!I\!\cdot\!R}
\newcommand{\CCU}{\text{CCU}}

% ========= 섹션 서식 =========
\titleformat{\section}{\LARGE\bfseries\color{darkblue}}{\thesection}{1em}{}
\titleformat{\subsection}{\normalsize\bfseries\color{darkblue}}{\thesubsection}{1em}{}
\titleformat{\subsubsection}{\normalsize\itshape}{\thesubsubsection}{1em}{}

% ========= 헤더/푸터 =========
\fancypagestyle{firstpage}{%
	\fancyhf{}%
	\renewcommand{\headrulewidth}{-5pt}%
	\renewcommand{\footrulewidth}{-5pt}%
	\fancyfoot[C]{%
		\begin{minipage}[t]{\textwidth}
			\centering
			\textcolor{gold}{\rule{\textwidth}{1.6pt}}\\[5pt]
			{\small \textsuperscript{1}Yakushev Research, \url{https://yuct.org/}} \hfill
			{\small YPSDC 프로토콜: \url{https://ypsdc.com/}}\\[-2pt]
			\textcolor{gold}{\rule{\textwidth}{1.6pt}}\\[5pt]
			\textbf{야쿠셰프 통일 조정 이론 2026} \hfill \thepage\\[10pt]
		\end{minipage}%
	}%
}

\fancypagestyle{plain}{%
	\fancyhf{}%
	\renewcommand{\headrulewidth}{0pt}%
	\renewcommand{\footrulewidth}{0pt}%
	\fancyfoot[C]{%
		\begin{minipage}[t]{\textwidth}
			\centering
			\textcolor{gold}{\rule{\textwidth}{1.6pt}}\\[4pt]
			\textbf{야쿠셰프 통일 조정 이론 2026} \hfill \thepage\\[4pt]
		\end{minipage}%
	}%
}

\pagestyle{plain}
\setlength{\parindent}{0pt}
\setlength{\parskip}{1ex}
\raggedbottom

\hypersetup{
	colorlinks=true,
	linkcolor=blue,
	citecolor=blue,
	urlcolor=blue,
}

% ========= 헤더 (YUCT 로고) – 가설/로드맵으로 변경 =========
\newcommand{\makeheader}{%
	\noindent
	\begin{minipage}{0.17\textwidth}
		\raggedright
		{\sffamily\bfseries\color{skyblue}\fontsize{46}{46}\selectfont YUCT}%
	\end{minipage}%
	\hspace{30pt}
	\begin{minipage}{0.32\textwidth}
		\raggedright
		{\sffamily\fontsize{11}{13}\selectfont 야쿠셰프\\ 통일\\ 조정 이론}%
	\end{minipage}%
	\hfill
	\begin{minipage}{0.38\textwidth}
		\raggedleft
		{\sffamily\bfseries 가설 / 로드맵}\\ 
		{\sffamily 버전 2.0 / 2026년 6월}\\ 
	\end{minipage}
	\par\vspace{5pt}
	\noindent\textcolor{gold}{\rule{\textwidth}{1.6pt}}
	\par\vspace{0pt}
}

\begin{document}
	\thispagestyle{firstpage}
	\makeheader
	
	\vspace{0cm}
	{\LARGE\bfseries 강한 CP 문제:\\[0.1cm] YUCT 안에서의 구체적 대수적 가설\\[0.2cm] \Large 이론적・실험적 검증을 위한 로드맵 \par}
	\vspace{0.2cm}
	
	{\large 알렉세이 V. 야쿠셰프\textsuperscript{1}} \\
	\vspace{0.0cm}
	
	{\color{darkblue}\textbf{\LARGE 초록}} \par
	{\large
		\noindent
		강한 CP 문제 — 왜 QCD의 CP를 깨는 매개변수 \(\theta\) 가 \(10^{-10}\) 보다 작은가 — 는 야쿠셰프 통일 조정 이론(YUCT) 안에서 표준 모형의 마지막 남은 수수께끼이다. 이 문서에서 우리는 필요한 억제를 자연스럽게 제공하는 구체적인 대수적 가설을 제시한다. 우리는 \(\theta\) 매개변수의 유효 조정 깊이가
		\[
		L_\theta = \frac{1}{2}\,L_W \;+\; \frac{S_{\mathrm{odd}}}{S_{\mathrm{even}}},
		\]
		로 주어진다고 제안한다. 여기서 \(L_W = 96\) 은 \(W\) 보손의 깊이이고 \(S_{\mathrm{odd}}/S_{\mathrm{even}} = 3/2\) 는 이론의 근본 상수이다. 이 공식은 자유 매개변수를 전혀 포함하지 않으며 \(L_\theta = 49.5\) 를 제공하고, 이는 \(\theta \sim (2/3)^{49.5} \approx 2.3 \times 10^{-10}\) 에 해당한다. 우리는 이 표현식의 물리적 동기를 논의하고, 그것을 YUCT의 더 넓은 대수적 틀 안에 포함시키며, 주로 중성자 전기 쌍극자 모멘트(nEDM)를 통해 이 가설을 반증할 수 있는 실험적 시험 계획의 개요를 서술한다. 이 제안은 미래의 데이터에 의한 반박이나 확인에 열려 있는 명료하고 수학적으로 정밀한 추측으로 제시된다.
	}
	
	\vspace{0.1cm}
	\noindent
	{\color{darkblue}\textbf{키워드:}} YUCT, 강한 CP 문제, 대수적 가설, 조정 깊이, 중성자 전기 쌍극자 모멘트, 열린 연구 프로그램.
	\vspace{0.4cm}
	
	\clearpage
	\tableofcontents
	\newpage
	
	\section{서론}
	\label{sec:intro}
	
	야쿠셰프 통일 조정 이론(YUCT)은 계층 문제, 우주 상수 문제, 페르미온 질량의 패턴 및 기초 물리학의 여러 주요 난제들에 대해 매개변수 없는 정량적 해결을 제공해 왔다~\cite{Y1,Y3,Y4}. 유일한 예외는 강한 CP 문제였다: 양자 색역학(QCD)의 CP-위반 \(\theta\)-매개변수가 부자연스러울 정도로 작아서 실험적 상한이 \(|\theta| \lesssim 10^{-10}\) 으로 제한된다는 것이다~\cite{nEDM2020}.
	
	YUCT 안에서 초기 시도들은 임의의 정수 \(k_\theta\) 를 도입하여 \(\theta\) 에 조정 깊이를 할당하는 것이었지만, 이는 구조적 예측이 아니라 피팅에 지나지 않아 만족스럽지 않았다. 이 노트에서 우리는 다른 접근법을 제시한다. 즉, QCD 진공의 질서 있는(페르미온) 섹터와 카오스적인(위상적) 섹터 사이의 상호 작용을 이용하는 것이다. YUCT는 질서와 카오스에 대한 통일된 대수적 기술 덕분에 이 상호 작용을 정량화할 수 있는 독특한 위치에 있다~\cite{Y2,Feig}.
	
	우리는 19차원 라그랑지안으로부터의 엄밀한 유도를 주장하지 않는다; 그러한 유도는 장기적인 목표이다. 대신에 우리는 YUCT에 이미 존재하는 숫자들로부터 자연스럽게 따라 나오고, 조정 가능한 매개변수를 전혀 포함하지 않으며, 중성자 전기 쌍극자 모멘트(nEDM)와 액시온 질량에 대해 반증 가능한 예측을 하는 날카로운 대수적 추측을 제시한다. 이 추측은 미래의 이론적·실험적 연구를 위한 명확한 목표로 제시된다.
	
	\section{YUCT에서의 질서와 카오스 상수들}
	\label{sec:oc}
	
	\subsection{질서}
	보편 오차 법칙
	\[
	\varepsilon = \kappa_c\,\alpha\,K_{\mathrm{eff}}^{-\beta},\qquad \beta = \frac{2}{3},\quad \kappa_c = \frac{1}{3},
	\]
	은 조정의 계층적 구조와 함께 닫힌 대수적 루프를 이끌어낸다:
	\[
	S_{\mathrm{odd}} = \frac{6}{5} = 1.2,\quad S_{\mathrm{even}} = \frac{4}{5} = 0.8,\quad 
	\frac{S_{\mathrm{odd}}}{S_{\mathrm{even}}} = \frac{3}{2} = \frac{1}{\beta}.
	\]
	모든 물리적 질량은 플랑크 질량과 조정 깊이 \(L\) 을 통해 표현된다:
	\[
	m \approx M_{\mathrm{Pl}}\left(\frac{2}{3}\right)^L.
	\]
	특히 \(W\) 보손의 질량은 \(L_W = 96\) 에 해당하는데, 이는 표준 모형(오른손잡이 중성미자 포함)의 바일 페르미온 자유도의 총 수이다~\cite{Y3}.
	
	\subsection{카오스}
	보편적인 파이겐바움 상수들
	\[
	\delta \approx 4.6692,\qquad \alpha \approx 2.5029,
	\]
	은 카오스로의 주기 배가 경로를 기술한다. YUCT에서 이들은 황금비 \(\varphi = (1+\sqrt{5})/2\) 와 원주율 \(\pi\) 을 통해 질서 상수들과 대수적으로 연결된다:
	\[
	2\alpha^2 \approx \pi\delta,\qquad
	\pi \approx S_{\mathrm{odd}}\varphi^2 \quad \Longrightarrow \quad 
	2\alpha^2 \approx (S_{\mathrm{odd}}\varphi^2)\,\delta .
	\]
	이 연결은 파이겐바움 상수들이 YUCT에 외래적인 것이 아니라 동일한 확장된 대수계에 속한다는 것을 보여준다~\cite{Y2}.
	
	\subsection{질서-카오스 공액으로서의 양자 진공}
	QCD 진공은 질서 있는 측면(쿼크 응축, 카이랄 대칭 깨짐)과 카오스적 측면(인스탄톤 요동, 위상적 감수율)이 공존하는 장이다. \(\theta\) 매개변수는 위상적 전하 밀도 \(G\widetilde{G}\) 와 결합하고, 두 섹터 모두로부터 기여를 받는다. 따라서 유효 깊이 \(L_\theta\) 가 질서 성분과 카오스 성분을 특징짓는 깊이들의 함수일 것이라고 기대하는 것은 자연스럽다.
	
	이전 연구에서 우리는 잠정적인 깊이 \(L_{\mathrm{ord}}\)(페르미온 섹터)와 \(L_{\mathrm{ch}}\)(위상적/카오스 섹터)를 도입했지만, 이들의 단순한 선형 결합은 필요한 억제 \(|\theta| \sim 10^{-10}\) 을 재현하지 못했다. 이것이 우리로 하여금 더 구조적인 결합을 찾도록 이끌었다.
	
	\section{\(L_\theta\) 에 대한 대수적 가설}
	\label{sec:hypothesis}
	
	\subsection{가설의 진술}
	
	우리는 \(\theta\) 매개변수의 유효 조정 깊이가 다음과 같은 단순한 공식으로 주어진다고 제안한다:
	\begin{equation}
		\boxed{L_\theta = \frac{1}{2}\,L_W \;+\; \frac{S_{\mathrm{odd}}}{S_{\mathrm{even}}}},
		\label{eq:Ltheta}
	\end{equation}
	여기서
	\begin{itemize}
		\item \(L_W = 96\) 은 \(W\) 보손의 깊이(전약 스케일의 기준선);
		\item \(S_{\mathrm{odd}}/S_{\mathrm{even}} = 3/2 = 1.5\) 는 세대 간 질량 인자를 지배하고 1차 대수적 루프에 나타나는 보편적 비율이다.
	\end{itemize}
	
	이 숫자들을 대입하면,
	\[
	L_\theta = \frac{96}{2} + 1.5 = 48 + 1.5 = 49.5
	\]
	를 얻는다.
	
	인스탄톤 계산에서 나올 수 있는 차수 1의 가능한 인자를 제외하면, 해당 \(\theta\) 매개변수는
	\[
	\theta \sim \left(\frac{2}{3}\right)^{L_\theta} = \left(\frac{2}{3}\right)^{49.5} \approx 2.3 \times 10^{-10}
	\]
	이다. 이 값은 실험적으로 허용된 범위 \(|\theta| \lesssim 10^{-10}\) 안에 정확히 들어온다~\cite{nEDM2020}.
	
	\subsection{물리적 해석}
	
	왜 이 특정한 결합이 나타나는가?
	\begin{enumerate}
		\item \textbf{인자 \(1/2\)}: 전약 스케일은 진공 기대값 \(v \approx 246\) GeV에 의해 정해지며, 이것은 기하학적 인자 \(\xi_W \approx 0.53\) 을 통해 기준 깊이 96으로부터 얻어진다. 깊이 48은 스케일 \(\sqrt{v M_{\mathrm{Pl}}}\) — 약한 스케일과 플랑크 스케일의 기하 평균 — 에 해당하며, 이는 페르미온과 보손의 진공 에너지 기여가 교차할 것으로 예상되는 바로 그 스케일이다. YUCT에서 이 스케일은 힉스 응축물의 조정 동역학에서 특별한 역할을 하는 것으로 알려져 있다(부록 \(\Lambda\)).
		\item \textbf{이동 \(S_{\mathrm{odd}}/S_{\mathrm{even}} = 3/2\)}: 비율 \(3/2 = 1/\beta\) 는 임의의 계층적 조정 체계에서 질서 있는(홀수 층) 상관관계와 무질서한(짝수 층) 상관관계 사이의 환원 불가능한 비대칭성을 정량화한다. 이것이 \(L_\theta\) 에 나타난다는 것은 CP를 위반하는 항이 진공의 질서 지배 섹터와 카오스 지배 섹터 사이의 작고 피할 수 없는 불균형에서 발생한다는 신호이다. 이 이동이 곱셈적이 아니라 덧셈적이라는 사실은 조정 네트워크의 이산적이고 층상 구조를 반영한다.
	\end{enumerate}
	
	따라서 식~(\ref{eq:Ltheta})은 다음과 같이 읽을 수 있다: \(\theta\) 의 깊이는 약한 깊이의 절반에 보편적인 질서-카오스 오프셋을 더한 것이다. 이 식은 자유 매개변수가 없으며, YUCT 안에서 이미 독립적으로 고정된 숫자들만을 사용한다.
	
	\subsection{직접적인 결과}
	
	\paragraph{중성자 전기 쌍극자 모멘트}
	nEDM은 \(d_n \approx 3.6 \times 10^{-16}\,\theta\; e\!\cdot\!\mathrm{cm}\) (강입자 불확실성 인자 포함)으로 \(\theta\) 와 관계된다~\cite{nEDM2020}. \(\theta \approx 2.3 \times 10^{-10}\) 를 취하면
	\[
	d_n \approx 8.3 \times 10^{-26}\; e\!\cdot\!\mathrm{cm}
	\]
	을 얻는다. 현재의 실험적 한계는 \(d_n < 1.8 \times 10^{-26}\; e\!\cdot\!\mathrm{cm}\) (90\% 신뢰도) 수준이다. PSI, SNS, TUCAN에서의 차세대 실험들은 \(\sim 10^{-28}\; e\!\cdot\!\mathrm{cm}\) 의 감도를 목표로 한다. 그 수준에서의 영(零) 결과는 이 가설을 배제할 것이고, 양성 검출은 극적인 확인이 될 것이다.
	
	\paragraph{액시온 물리}
	만약 강한 CP 문제가 페체이-퀸 메커니즘으로 해결된다면, 액시온 질량 \(m_a\) 는 \(\theta\) 및 액시온 붕괴 상수 \(f_a\) 와 관계된다. 식~(\ref{eq:Ltheta})은 동일한 조정 깊이 논리를 통해 \(f_a\)(그리고 따라서 \(m_a\))를 고정한다. 이것은 ADMX 및 다른 액시온 탐색에 상보적인 실험 목표를 제공한다.
	
	\section{검증을 위한 로드맵}
	\label{sec:roadmap}
	
	이 가설은 수학적으로 정밀하고 반증 가능하다. 우리는 다음과 같은 프로그램을 제안한다.
	
	\subsection{이론적 과제들}
	\begin{enumerate}
		\item \textbf{19차원 라그랑지안으로부터의 유도} YUCT 19차원 모이론의 유효 작용을 표준 QCD \(G\widetilde{G}\) 항으로 축소해야 한다. 올바른 진공 상태가 확인될 때, 축소 과정에서의 계수가 자연스럽게 깊이 \(L_\theta = 49.5\) 를 산출한다는 것을 보이는 것이 목표이다. 이것은 도전적이지만 잘 정의된 수학적 문제이다.
		\item \textbf{격자 게이지 이론} 기존의 격자 QCD 계산은 위상적 감수율과 진공 에너지의 \(\theta\) 의존성을 계산할 수 있다. YUCT에서 영감을 받은 분석은 식~(\ref{eq:Ltheta})의 대수적 구조가 결합 상수에 따른 격자 결과의 축척 행동이나 위상적 전하 클러스터의 분포에 반영되는지 검증해야 한다.
		\item \textbf{nEDM 예측의 정밀화} \(\theta\) 와 \(d_n\) 을 연결하는 강입자 행렬 요소는 카이랄 섭동 이론과 격자 QCD로부터 알려져 있다. 이 행렬 요소들의 전용 계산과 YUCT의 \(\theta\) 값을 결합하면, 차세대 nEDM 한계와 직접 비교할 수 있는 날카로운 예측을 얻을 수 있을 것이다.
	\end{enumerate}
	
	\subsection{실험적 과제들}
	\begin{enumerate}
		\item \textbf{nEDM 실험 (PSI, SNS, TUCAN)} 측정된 nEDM이 \(10^{-26}\;e\!\cdot\!\mathrm{cm}\) 이상이면 가설과 부합할 것이고, \(10^{-28}\;e\!\cdot\!\mathrm{cm}\) 미만의 한계는 가설을 배제할 것이다.
		\item \textbf{액시온 탐색 (ADMX, MADMAX, IAXO)} 만약 페체이-퀸 메커니즘이 실현된다면, \(L_\theta\) 로부터 \(f_a\) 를 평가하여 YUCT가 예측하는 액시온 질량을 계산할 수 있을 것이다. 해당 질량 창에서의 전용 탐색은 가설을 시험할 것이다.
		\item \textbf{우주론적 흔적} 질서-카오스 간섭이 근본적이라면, 초기 우주의 상전이에도 영향을 주어 우주 마이크로파 배경이나 암흑 물질 분포에 특정한 흔적을 남길 수 있다. 이것들은 Planck, Euclid 및 미래의 CMB-S4 데이터로 제약될 수 있다.
	\end{enumerate}
	
	\section{협력에의 초대}
	\label{sec:invite}
	
	YUCT 프레임워크는 완전한 라그랑지안, 모든 부록, 그리고 방대한 조정 깊이 데이터베이스를 포함하여 Zenodo에 공개되어 있다~\cite{Y1,Y2}. 우리는 물리학자, 수학자, 전산 과학자들이 여기 제시된 가설을 유도하거나, 시험하거나, 반증하는 노력에 동참하기를 초대한다. 저자(알렉세이 야쿠셰프, \href{mailto:alexey@yakushev.eu}{alexey@yakushev.eu})는 서신, 제안 및 협력 제의를 환영한다.
	
	\section{결론}
	\label{sec:conclusion}
	
	우리는 야쿠셰프 통일 조정 이론 안에서 QCD \(\theta\) 매개변수의 유효 깊이에 대해, 매개변수가 없는 구체적인 대수적 공식을 제안했다:
	\[
	L_\theta = \frac{1}{2}L_W + \frac{S_{\mathrm{odd}}}{S_{\mathrm{even}}} = 49.5,
	\]
	이것은 \(\theta \approx 2.3 \times 10^{-10}\) 을 산출한다. 이 가설은 수학적으로 단순하고, 물리적으로 동기가 있으며, 다가오는 nEDM 및 액시온 실험에 의해 직접적으로 반증 가능하다. 이것은 임의의 정수 할당에 의존했던 이전의 시도들을 대체하며, YUCT 내에서 강한 CP 문제의 완전한 해결로 가는 명확한 길을 연다. 가설이 실험적 검증을 견뎌낼지는 아직 알 수 없지만, 그 가치는 미래 연구를 이끌 정밀하고 시험 가능한 추측을 제공한다는 데 있다.
	
	\vspace{1cm}
	\noindent\textbf{날짜:} 2026년 6월\\
	\textbf{버전:} 2.0\\
	\textbf{상태:} 구체적 대수적 예측을 수반하는 가설; 실험적 시험과 이론적 유도에 열려 있음.
	
	\newpage
	\begin{thebibliography}{99}
		\bibitem{Y1} A.V. Yakushev, \textit{Yakushev Unified Coordination Theory (YUCT)}, Zenodo, 2026. DOI:10.5281/zenodo.18444599.
		\bibitem{Y2} A.V. Yakushev, Appendix AF: The Algebraic Framework of Reality in YUCT, 2026.
		\bibitem{Y3} A.V. Yakushev, Appendix \(\Lambda\): Resolution of the Hierarchy and Cosmological Constant Problems, 2026.
		\bibitem{Y4} A.V. Yakushev, Appendix J: Complete Derivation of Standard Model Flavor Parameters, 2026.
		\bibitem{Feig} M.J. Feigenbaum, \textit{J. Stat. Phys.} \textbf{19}, 25 (1978).
		\bibitem{nEDM2020} C. Abel et al. (nEDM Collaboration), \textit{Phys. Rev. Lett.} \textbf{124}, 081803 (2020).
	\end{thebibliography}
	
\end{document}